\section{Hand of Straights}
\label{sec:heap-hand-of-straights}

\problemstatement{We are given a hand of $n$ cards, $\textit{hand}[0..n-1]$
(integers, possibly with repeats), and an integer $\textit{groupSize}$.
Determine whether the hand can be rearranged into groups, each of exactly
$\textit{groupSize}$ cards with \textbf{consecutive} values (e.g.\
$[3,4,5]$ is a valid group of size $3$; $[3,4,6]$ is not). Return
\texttt{true} or \texttt{false}.}

\begin{patternbox}
\emph{``Always form the next group starting from the smallest remaining
card''} is the greedy insight, and \emph{``repeatedly find the current
smallest remaining distinct value''} is exactly what a min-heap (over
distinct values) is built for. Whenever a problem requires repeatedly
peeling off the globally smallest (or largest) remaining item from a
\emph{changing} collection -- not just once, but again and again as the
collection shrinks -- that repeated ``smallest of what's left'' query is
the heap keyword to watch for, distinguishing this from a single-pass
greedy that only needs one sorted scan.
\end{patternbox}

\subsection*{Understanding the problem carefully}

First, if $n$ is not divisible by $\textit{groupSize}$, the answer is
immediately \texttt{false} -- the cards cannot be partitioned evenly.
Otherwise, count the occurrences of each distinct card value (a hash map),
and observe a key greedy fact: whichever card value is currently the
\textbf{smallest} among all cards not yet placed into a group \emph{must}
be the starting card of some group -- it cannot be the second or later
card of any group, since no smaller card remains to precede it. So we are
forced to start a new group at the current minimum, every time.

\begin{ideabox}[Greedy Choice, Powered by a Min-Heap]
Maintain a min-heap of the \emph{distinct} card values currently present
in the hand (with their counts tracked separately in a hash map). Repeatedly: peek
the smallest remaining value $v$; this must start a new group. Check that
$v, v+1, \dots, v+\textit{groupSize}-1$ are all available (count $>0$ in
the hash map); if any is missing, return \texttt{false}. Otherwise,
decrement the count of each of these $\textit{groupSize}$ values; whenever
a value's count reaches $0$, remove it from consideration (pop it from the
heap, or simply leave it and skip over exhausted values when they resurface
at the top). Repeat until all cards are consumed.
\end{ideabox}

\subsection*{Why always starting at the current minimum is correct}

This is another exchange-argument greedy: suppose, for contradiction, that
some valid grouping does \emph{not} use the smallest remaining card $v$ as
the start of a group, but instead places it later in some group starting
at $v' < v$ -- but $v$ \emph{is} the minimum remaining, so no $v' < v$ can
exist among the remaining cards. Therefore $v$ cannot appear anywhere
except as the start of a group (it cannot be the second card of a group
starting at $v-1$, because $v-1$ is not present among the remaining
cards by assumption that $v$ is the minimum). So every valid grouping is
forced to start a group at $v$, and the greedy choice has no alternative
to consider in the first place -- it is not just \emph{a} good choice, it
is the \emph{only} choice consistent with feasibility.

\subsection*{Algorithm}

\begin{enumerate}
  \item If $n \bmod \textit{groupSize} \ne 0$: return \texttt{false}.
  \item Build a frequency map \textit{count} of card values; build a
        min-heap of the distinct values present.
  \item While the heap is not empty:
  \begin{enumerate}
    \item Pop values from the top of the heap whose count has reached
          $0$ (already fully used) and discard them.
    \item Peek the new minimum $v$ (if the heap is now empty, stop --
          all cards are used).
    \item For $j \gets 0$ to $\textit{groupSize}-1$: if
          $\textit{count}[v+j] = 0$, return \texttt{false}; otherwise
          decrement $\textit{count}[v+j]$.
  \end{enumerate}
  \item Return \texttt{true}.
\end{enumerate}

\begin{algorithm}[H]
\caption{Hand of Straights}
\begin{algorithmic}[1]
\Procedure{IsNStraightHand}{\textit{hand}, \textit{groupSize}}
  \If{$n \bmod \textit{groupSize} \ne 0$} \Return \textbf{false} \EndIf
  \State build frequency map \textit{count}; build min-heap of distinct values
  \While{\textit{heap} is not empty}
    \While{\textit{heap} not empty \textbf{and} $\textit{count}[\text{top}] = 0$}
      \State pop \textit{heap}
    \EndWhile
    \If{\textit{heap} is empty} \textbf{break} \EndIf
    \State $v \gets$ top of \textit{heap}
    \For{$j \gets 0$ to $\textit{groupSize}-1$}
      \If{$\textit{count}[v+j] = 0$} \Return \textbf{false} \EndIf
      \State $\textit{count}[v+j] \gets \textit{count}[v+j] - 1$
    \EndFor
  \EndWhile
  \State \Return \textbf{true}
\EndProcedure
\end{algorithmic}
\end{algorithm}

\subsection*{Dry run}

\dryrun{Take $\textit{hand} = [1,2,3,6,2,3,4,7,8]$,
$\textit{groupSize} = 3$.}

$n=9$, $9 \bmod 3 = 0$. Counts: $1{:}1, 2{:}2, 3{:}2, 4{:}1, 6{:}1, 7{:}1,
8{:}1$. Min-heap of distinct values: $\{1,2,3,4,6,7,8\}$.

\begin{itemize}
  \item Min is $1$. Group $[1,2,3]$: all counts $>0$. Decrement:
        $1{:}0, 2{:}1, 3{:}1$.
  \item Top of heap is $1$, but $\textit{count}[1]=0$ now; pop it. New
        min is $2$. Group $[2,3,4]$: all counts $>0$. Decrement:
        $2{:}0, 3{:}0, 4{:}0$.
  \item Pop $2,3,4$ as their counts reach $0$ at the top. New min is $6$.
        Group $[6,7,8]$: all counts $>0$ ($1$ each). Decrement to $0$
        each.
  \item Heap empties out as exhausted values are popped.
\end{itemize}

All cards consumed without failure, so the answer is \texttt{true}: the
hand splits into $[1,2,3]$, $[2,3,4]$, $[6,7,8]$.

\subsection*{C++ implementation}

\begin{lstlisting}
#include <bits/stdc++.h>
using namespace std;

bool isNStraightHand(vector<int>& hand, int groupSize) {
    int n = hand.size();
    if (n % groupSize != 0) return false;

    unordered_map<int,int> count;
    for (int x : hand) count[x]++;

    priority_queue<int, vector<int>, greater<int>> minHeap;
    for (auto& [val, c] : count) minHeap.push(val);

    while (!minHeap.empty()) {
        while (!minHeap.empty() && count[minHeap.top()] == 0) {
            minHeap.pop();
        }
        if (minHeap.empty()) break;

        int v = minHeap.top();
        for (int j = 0; j < groupSize; j++) {
            if (count[v + j] == 0) return false;
            count[v + j]--;
        }
    }
    return true;
}

int main() {
    vector<int> hand = {1, 2, 3, 6, 2, 3, 4, 7, 8};
    cout << (isNStraightHand(hand, 3) ? "true" : "false") << endl;
    return 0;
}
\end{lstlisting}

\begin{complexitybox}
\textbf{Time Complexity.} Building the frequency map and heap costs
$O(n \log n)$ in the worst case (up to $n$ distinct values, each pushed
once). The main loop performs at most $n/\textit{groupSize}$ group
formations, each doing $O(\textit{groupSize})$ work plus
$O(\log n)$ amortized heap-popping of exhausted values, giving an overall
time complexity of
\[
  O(n \log n).
\]
\textbf{Space Complexity.} $O(n)$ for the frequency map and heap in the
worst case (all distinct values).
\end{complexitybox}

\begin{takeawaybox}
When a greedy choice is not just \emph{good} but \emph{forced} (the
smallest remaining element has no alternative role to play), a min-heap
is the natural data structure to repeatedly surface that forced choice
efficiently, especially when -- unlike a single sorted pass -- elements
get consumed and removed from consideration in a pattern that does not
simply move left to right.
\end{takeawaybox}
